\documentclass[aspectratio=169,12pt]{beamer}

% ============================================================
% PRESENTATION ORALE - HandControlPC
% Structure conforme a l'enonce NFP101 (section 5 : support oral)
% Compilation : tectonic HandControlPC_Presentation.tex
% Conversion PPT : ouvrir le PDF dans PowerPoint/Keynote ou via un convertisseur
% ============================================================

% --- Encodage et langue ---
\usepackage[utf8]{inputenc}
\usepackage[T1]{fontenc}
\usepackage[french]{babel}

% --- Theme ---
\usetheme{Madrid}
\usecolortheme{whale}
\useinnertheme{circles}
\setbeamertemplate{navigation symbols}{}
\setbeamertemplate{caption}[numbered]

% --- Couleurs personnalisees ---
\usepackage{xcolor}
\definecolor{cnamblue}{RGB}{25,60,120}
\definecolor{steelblue}{RGB}{70,130,180}
\definecolor{accent}{RGB}{200,100,20}
\setbeamercolor{structure}{fg=cnamblue}
\setbeamercolor{frametitle}{bg=cnamblue,fg=white}
\setbeamercolor{title}{fg=white}

% --- Packages ---
\usepackage{graphicx}
\usepackage{booktabs}
\usepackage{tabularx}
\usepackage{tikz}
\usetikzlibrary{shapes.geometric, arrows.meta, positioning}
\usepackage{listings}
\usepackage{amsmath}

% --- Style code ---
\lstset{
  basicstyle=\ttfamily\scriptsize,
  backgroundcolor=\color{gray!10},
  frame=single,
  rulecolor=\color{gray!40},
  breaklines=true,
  showstringspaces=false,
  commentstyle=\color{gray},
  keywordstyle=\color{cnamblue}\bfseries,
  stringstyle=\color{green!50!black},
  language=Python,
}

% --- Barre de progression simple en pied de page (deja dans Madrid) ---

% --- Infos titre ---
\title[HandControlPC]{HandControlPC}
\subtitle{Contr\^ole gestuel du PC par webcam}
\author[M.A. Sobhi]{Mohamed Amine Sobhi}
\institute[Cnam]{Cnam Occitanie --- NFP101 Programmation Orient\'ee Objet\\Professeur~: Adrien Escourrou}
\date{Mai 2026}

% ============================================================
\begin{document}

% --- DIAPO TITRE ---
\begin{frame}[plain]
  \titlepage
\end{frame}

% --- SOMMAIRE ---
\begin{frame}{Plan de la pr\'esentation}
  \tableofcontents
\end{frame}

% ============================================================
\section{Introduction, contexte et objectif}
% ============================================================

\begin{frame}{Contexte~: pourquoi HandControlPC~?}
  \begin{columns}
    \begin{column}{0.58\textwidth}
      \textbf{Le constat}
      \begin{itemize}
        \item Clavier et souris imposent un contact physique permanent.
        \item Certains contextes y sont mal adapt\'es~: pr\'esentation, accessibilit\'e, mains occup\'ees.
      \end{itemize}
      \vspace{0.4cm}
      \textbf{L'id\'ee}
      \begin{itemize}
        \item Utiliser la \textbf{webcam} d\'ej\`a pr\'esente sur chaque PC.
        \item Traduire des \textbf{gestes de la main} en actions syst\`eme.
        \item Aucun mat\'eriel suppl\'ementaire.
      \end{itemize}
    \end{column}
    \begin{column}{0.42\textwidth}
      \begin{block}{Public cible}
        \begin{itemize}
          \item Utilisateurs en pr\'esentation
          \item Accessibilit\'e (handicap moteur)
          \item Curieux de la vision par ordinateur
        \end{itemize}
      \end{block}
    \end{column}
  \end{columns}
\end{frame}

\begin{frame}{Objectif du projet}
  \begin{block}{Fonctionnalit\'e centrale}
    Capter les gestes de la main en temps r\'eel et les convertir en actions concr\`etes~:
    \textbf{volume}, \textbf{capture d'\'ecran}, \textbf{zoom}, et \textbf{reconnaissance de signes}.
  \end{block}
  \vspace{0.3cm}
  \begin{columns}
    \begin{column}{0.5\textwidth}
      \textbf{Objectifs techniques}
      \begin{itemize}
        \item Application \emph{enti\`erement finie} (pas un prototype)
        \item Architecture orient\'ee objet propre
        \item Configuration externe
        \item Tests automatis\'es
      \end{itemize}
    \end{column}
    \begin{column}{0.5\textwidth}
      \textbf{Contraintes}
      \begin{itemize}
        \item Temps r\'eel ($>$ 30 FPS)
        \item Robustesse aux faux positifs
        \item Code lisible et extensible
      \end{itemize}
    \end{column}
  \end{columns}
\end{frame}

% ============================================================
\section{Les \'etapes du projet}
% ============================================================

\begin{frame}{Les diverses \'etapes du projet}
  \begin{center}
  \resizebox{\textwidth}{!}{%
  \begin{tikzpicture}[
    node distance=0.35cm,
    nbox/.style={rectangle, draw=cnamblue, fill=cnamblue!10, rounded corners=4pt,
                 minimum width=2.5cm, minimum height=1.1cm, align=center, font=\small},
    arrow/.style={-{Latex[length=3mm]}, very thick, draw=cnamblue}
  ]
    \node[nbox] (s1) {1. Prototype\\d\'etection main};
    \node[nbox, right=of s1] (s2) {2. 3 modes\\(vol/screen/zoom)};
    \node[nbox, right=of s2] (s3) {3. Refactor\\POO (h\'eritage)};
    \node[nbox, right=of s3] (s4) {4. Config JSON\\+ tests};
    \node[nbox, right=of s4] (s5) {5. Signes\\BONJOUR};
    \draw[arrow] (s1)--(s2);
    \draw[arrow] (s2)--(s3);
    \draw[arrow] (s3)--(s4);
    \draw[arrow] (s4)--(s5);
  \end{tikzpicture}%
  }
  \end{center}
  \vspace{0.5cm}
  \begin{itemize}
    \item D\'emarche \textbf{incr\'ementale}~: chaque \'etape ajoute une fonctionnalit\'e d\'emontrable.
    \item Suivi par \textbf{commits Git r\'eguliers} (preuve d'avancement).
    \item Refactorisation POO \emph{apr\`es} avoir valid\'e le fonctionnel.
  \end{itemize}
\end{frame}

% ============================================================
\section{D\'emonstration et r\'esultat}
% ============================================================

\begin{frame}{D\'emonstration --- les 4 modes}
  \begin{columns}
    \begin{column}{0.5\textwidth}
      \begin{block}{Modes de contr\^ole}
        \begin{enumerate}
          \item \textbf{Volume}~: pincement pouce-index
          \item \textbf{Screenshot}~: 3 doigts maintenus
          \item \textbf{Zoom}~: \'ecartement de 2 mains
          \item \textbf{Signes}~: \'epeler BONJOUR
        \end{enumerate}
      \end{block}
    \end{column}
    \begin{column}{0.5\textwidth}
      \begin{block}{R\'esultats observables}
        \begin{itemize}
          \item Volume syst\`eme modifi\'e en direct
          \item Fichier PNG horodat\'e + son
          \item Flux vid\'eo zoom\'e en temps r\'eel
          \item Synth\`ese vocale \og Bonjour \fg
        \end{itemize}
      \end{block}
    \end{column}
  \end{columns}
  \vspace{0.3cm}
  \begin{center}
    \fbox{\parbox{0.85\textwidth}{\centering
      \textbf{[ D\'emonstration en direct --- $\sim$3 minutes ]}\\
      \small Lancer \texttt{python main.py} puis \texttt{python test\_sign\_recognition.py}
    }}
  \end{center}
\end{frame}

\begin{frame}{Le geste \og BONJOUR \fg{} en pratique}
  \begin{center}
  \small
  \begin{tabular}{cccccccc}
    \textbf{B} & \textbf{O} & \textbf{N} & \textbf{J} & \textbf{O} & \textbf{U} & \textbf{R} \\
    \midrule
    4 doigts & 3 du milieu & pouce+2 & shaka & 3 du milieu & V & index \\
  \end{tabular}
  \end{center}
  \vspace{0.4cm}
  \begin{itemize}
    \item Chaque lettre valide apr\`es \textbf{20 frames stables} ($\sim$0,7\,s).
    \item Affichage progressif~: \textcolor{green!60!black}{vert} = valid\'e, \textcolor{orange}{jaune} = attendu, gris = \`a venir.
    \item Une erreur de s\'equence remet le compteur \`a z\'ero.
    \item S\'equence compl\`ete $\rightarrow$ synth\`ese vocale + message de succ\`es.
  \end{itemize}
\end{frame}

% ============================================================
\section{Architecture technique}
% ============================================================

\begin{frame}{Architecture~: h\'eritage et polymorphisme}
  \begin{columns}[T]
    \begin{column}{0.48\textwidth}
      \begin{center}
      \resizebox{\linewidth}{!}{%
      \begin{tikzpicture}[
        node distance=0.9cm and 0.35cm,
        cbox/.style={rectangle, draw=cnamblue, fill=cnamblue!8, rounded corners=3pt,
                     minimum width=2.6cm, minimum height=0.8cm, align=center, font=\scriptsize\ttfamily},
        abox/.style={cbox, fill=steelblue!25, draw=steelblue, font=\scriptsize\ttfamily\bfseries},
        ar/.style={-{Latex[length=2.5mm]}, thick, draw=cnamblue}
      ]
        \node[abox] (ga) {GestureAction (ABC)};
        \node[cbox, below=of ga] (st) {ScreenshotTaker};
        \node[cbox, left=of st] (vc) {VolumeController};
        \node[cbox, right=of st] (zm) {ZoomManager};
        \draw[ar] (vc)--(ga);
        \draw[ar] (st)--(ga);
        \draw[ar] (zm)--(ga);
      \end{tikzpicture}%
      }
      \end{center}
    \end{column}
    \begin{column}{0.52\textwidth}
      \begin{block}{Choix cl\'es}
        \begin{itemize}
          \item \texttt{process()} \textbf{abstraite} : chaque mode sa logique.
          \item \texttt{render\_hud()} \textbf{h\'erit\'ee} : affichage commun.
          \item \textbf{Open/Closed} : ajouter un mode = une sous-classe.
        \end{itemize}
      \end{block}
    \end{column}
  \end{columns}
  \vspace{0.3cm}
  \begin{exampleblock}{Polymorphisme}
    \texttt{main.py} manipule les actions via l'interface commune \texttt{GestureAction}.
  \end{exampleblock}
\end{frame}

\begin{frame}{Pipeline temps r\'eel et stack technique}
  \begin{center}
  \resizebox{\textwidth}{!}{%
  \begin{tikzpicture}[
    node distance=0.3cm,
    nbox/.style={rectangle, draw=cnamblue, fill=cnamblue!8, rounded corners=3pt,
                 minimum width=1.9cm, minimum height=0.7cm, align=center, font=\scriptsize},
    arrow/.style={-{Latex[length=2mm]}, thick, draw=cnamblue}
  ]
    \node[nbox] (cam)  {Webcam\\OpenCV};
    \node[nbox, right=of cam] (mp)  {MediaPipe};
    \node[nbox, right=of mp]  (lm)  {21 points};
    \node[nbox, right=of lm]  (fp)  {fingers\_up};
    \node[nbox, right=of fp]  (md)  {Mode actif};
    \node[nbox, right=of md]  (sys) {Action OS};
    \draw[arrow] (cam)--(mp); \draw[arrow] (mp)--(lm); \draw[arrow] (lm)--(fp);
    \draw[arrow] (fp)--(md);  \draw[arrow] (md)--(sys);
  \end{tikzpicture}%
  }
  \end{center}
  \vspace{0.4cm}
  \begin{columns}
    \begin{column}{0.5\textwidth}
      \begin{block}{Biblioth\`eques}
        \begin{itemize}
          \item \textbf{MediaPipe}~: 21 landmarks/main
          \item \textbf{OpenCV}~: vid\'eo + HUD
          \item \textbf{NumPy}~: interp / clip
          \item \textbf{MSS}~: capture d'\'ecran
        \end{itemize}
      \end{block}
    \end{column}
    \begin{column}{0.5\textwidth}
      \begin{block}{Configuration externe}
        \texttt{config.json} centralise tous les param\`etres~: r\'esolution,
        seuils, dossiers. \\Chargement avec \texttt{deepcopy} + fallback.
      \end{block}
    \end{column}
  \end{columns}
\end{frame}

% ============================================================
\section{Validation et tests}
% ============================================================

\begin{frame}{Validation, tests et robustesse}
  \begin{columns}
    \begin{column}{0.52\textwidth}
      \begin{block}{29 tests unitaires (pytest)}
        \begin{itemize}
          \item \textbf{Sans webcam} : d\'ependances mock\'ees.
          \item Signes, s\'equence BONJOUR, volume, zoom, config.
          \item \texttt{python -m pytest tests/ -v}
        \end{itemize}
      \end{block}
      \begin{exampleblock}{R\'esultat}
        \textbf{29 passed} en $\sim$1,5\,s
      \end{exampleblock}
    \end{column}
    \begin{column}{0.48\textwidth}
      \begin{block}{Robustesse \og en direct \fg}
        \begin{itemize}
          \item \textbf{Cooldown} entre actions
          \item \textbf{Historique glissant} (anti faux positif)
          \item Seuil de confiance $\geq 0{,}6$
          \item \texttt{try/except} sur chaque frame
        \end{itemize}
      \end{block}
    \end{column}
  \end{columns}
\end{frame}

% ============================================================
\section{Discussion, limites et IA}
% ============================================================

\begin{frame}{Discussion --- limites et axes d'am\'elioration}
  \begin{columns}
    \begin{column}{0.5\textwidth}
      \begin{block}{Limites}
        \begin{itemize}
          \item macOS uniquement (osascript, afplay)
          \item Sensible \`a la luminosit\'e
          \item Signes \emph{statiques} seulement
          \item Mono-utilisateur
        \end{itemize}
      \end{block}
    \end{column}
    \begin{column}{0.5\textwidth}
      \begin{block}{Pistes}
        \begin{itemize}
          \item Contr\^ole souris par geste
          \item Mode apprentissage de gestes
          \item Portage Windows/Linux
          \item Logs persistants
        \end{itemize}
      \end{block}
    \end{column}
  \end{columns}
\end{frame}

\begin{frame}{Utilisation de l'IA --- transparence}
  \begin{block}{IA d\'eclar\'ees}
    Code assist\'e encadr\'e par \texttt{\# CODE IA (Claude)} dans les fichiers sources concern\'es.
  \end{block}
  \vspace{0.15cm}
  \begin{columns}[T]
    \begin{column}{0.5\textwidth}
      \textbf{O\`u l'IA est marqu\'ee}
      \begin{itemize}
        \item \textbf{Claude} --- \texttt{volume\_control.py}\\(classe + \texttt{\_get/\_set\_volume})
        \item \textbf{Claude} --- \texttt{zoom\_control.py}\\(m\'ethode \texttt{adjust()})
        \item \textbf{GitHub Copilot} --- \texttt{main.py}\\(refactorisation l\'eg\`ere)
      \end{itemize}
    \end{column}
    \begin{column}{0.5\textwidth}
      \textbf{Ce que j'ai ma\^itris\'e}
      \begin{itemize}
        \item Logique de d\'etection (\texttt{fingers\_up})
        \item Choix it\'eratif des patterns
        \item Int\'egration MediaPipe/OpenCV
        \item Architecture POO et tests
      \end{itemize}
    \end{column}
  \end{columns}
  \vspace{0.15cm}
  \begin{center}
    \small Je peux expliquer chaque passage assist\'e devant le jury.
  \end{center}
\end{frame}

% ============================================================
\section{Bilan et r\'ef\'erences}
% ============================================================

\begin{frame}{Ce que ce projet m'a apport\'e}
  \begin{columns}
    \begin{column}{0.5\textwidth}
      \begin{block}{Comp\'etences acquises}
        \begin{itemize}
          \item Conception \textbf{POO} concr\`ete (ABC, h\'eritage, polymorphisme)
          \item Vision par ordinateur temps r\'eel
          \item Tests automatis\'es et mocking
          \item Gestion de configuration
        \end{itemize}
      \end{block}
    \end{column}
    \begin{column}{0.5\textwidth}
      \begin{block}{Applicable en entreprise~?}
        \textbf{Oui}~:
        \begin{itemize}
          \item Architecture extensible (OCP)
          \item Tests $\rightarrow$ int\'egration continue
          \item Config externe $\rightarrow$ d\'eploiement
          \item Tra\c{c}abilit\'e (Git, doc, IA d\'eclar\'ee)
        \end{itemize}
      \end{block}
    \end{column}
  \end{columns}
\end{frame}

\begin{frame}{Ce que j'aurais chang\'e dans l'\'enonc\'e}
  \begin{itemize}
    \item[\textcolor{green!60!black}{$+$}] Sujet libre~: \textbf{tr\`es motivant}, permet de choisir un domaine qui passionne.
    \item[\textcolor{green!60!black}{$+$}] Exigences claires (POO, tests, doc)~: bon cadre p\'edagogique.
    \vspace{0.3cm}
    \item[\textcolor{accent}{$\sim$}] Pr\'eciser un \textbf{ordre de grandeur} attendu (nb de classes, lignes) pour cadrer la port\'ee.
    \item[\textcolor{accent}{$\sim$}] Proposer une \textbf{liste d'exemples} de sujets pour ceux qui h\'esitent.
    \item[\textcolor{accent}{$\sim$}] Donner un \textbf{barème d\'etaill\'e des bonus} pour mieux prioriser.
  \end{itemize}
\end{frame}

\begin{frame}{Bibliographie et r\'ef\'erences}
  \begin{block}{Biblioth\`eques et documentation}
    \begin{itemize}
      \item \textbf{MediaPipe} (Google) --- d\'etection de mains
      \item \textbf{OpenCV} --- vision par ordinateur
      \item \textbf{Python ABC} --- classes abstraites
      \item \textbf{pytest} --- tests unitaires
      \item \textbf{MSS} --- capture d'\'ecran
    \end{itemize}
  \end{block}
  \begin{block}{Intelligence artificielle (d\'eclar\'ee)}
    \begin{itemize}
      \item \textbf{Claude Code} (Anthropic) --- \texttt{volume\_control.py}, \texttt{zoom\_control.py}
      \item \textbf{GitHub Copilot} --- refactorisation l\'eg\`ere de \texttt{main.py}
    \end{itemize}
  \end{block}
\end{frame}

% --- DIAPO FINALE ---
\begin{frame}[plain]
  \begin{center}
    \vfill
    {\Huge \textbf{\textcolor{cnamblue}{Merci de votre attention}}}\\[0.6cm]
    {\Large Questions \& r\'eponses}\\[1cm]
    {\large HandControlPC --- Mohamed Amine Sobhi}\\[0.2cm]
    {\normalsize Cnam NFP101 --- Mai 2026}
    \vfill
  \end{center}
\end{frame}

\end{document}
